%----------HEADING-----------------
\begin{tabular*}{\textwidth}{l@{\extracolsep{\fill}}r}
  \textbf{\href{https://www.linkedin.com/in/yufeng-xia-12648a25b}{\Large Yufeng (Alex) Xia}} & Email: \href{mailto:yufeng.xia@ed.ac.uk}{yufeng.xia@ed.ac.uk} \\
  \href{https://www.linkedin.com/in/yufeng-xia-12648a25b}{linkedin.com/in/yufeng-xia-12648a25b} & Mobile: \href{tel:+447918155093}{(+44) 7918155093} \\
  Edinburgh, UK & \href{https://xyf2002.github.io/}{xyf2002.github.io} \\
\end{tabular*}

\vspace{-2pt}

%-----------PROFILE-----------------
\section{Profile}
  \begin{itemize}[leftmargin=*]
    \item\small{MSc by Research student at the Institute for Computing Systems Architecture (ICSA), University of Edinburgh, supervised by Prof.\ Mahesh K.\ Marina. My interest lies at the intersection of \textbf{wireless systems, systems for AI, and large-scale emulation}. My work spans \textbf{AI-RAN}, where live radio processing and data-driven workloads share the same infrastructure under real-time constraints, as well as \textbf{emulation and digital twin systems}, the platforms on which RAN and edge-AI algorithms are developed and validated at scale.
    \par\vspace{3pt}
    I build the layer between the radio and the model, making live network behaviour something models can be trained and validated against, and idle radio compute something they can run on. These systems have led to publications at top-tier venues including ACM MobiCom, and what interests me most now is carrying this line of work into next-generation radio access and edge AI infrastructure. \textbf{Available full-time from September 2026.}}
  \end{itemize}
